% TOP_PROBLEM: {"id": 1423, "title": "Snake-in-the-Box Problem", "category_id": 3, "category": {"id": 3, "name": "graph_theory", "display_name": "Graph Theory", "description": "Problems involving graphs, networks, and their properties.", "slug": "graph-theory", "order_index": 3, "created_at": "2026-07-31T15:26:25.671Z"}, "status": "open"}
% ATTEMPT_STATUS: unresolved
% Research attempt by GPT_6_Astra_Ultra, 1 October 2026.
% Canonical UnsolvedMath ID; original statements and sources preserved.
% FIRST_POSTED: 2026-10-01
% Imported from: unsolvedmath_additional_500.tex; UnsolvedMath ID 1423
% !TeX program = lualatex
% Compile twice: lualatex 1423.tex
% Self-contained: no external bibliography, images, or \input files.
% Numeric section labels and visible numbers are original UnsolvedMath IDs.
\documentclass[11pt,a4paper]{article}
\usepackage[margin=24mm,headheight=14pt]{geometry}
\usepackage{fontspec}
\setmainfont{Latin Modern Roman}
\setsansfont{Latin Modern Sans}
\setmonofont{Latin Modern Mono}
\usepackage{amsmath,amssymb,mathtools}
\usepackage{unicode-math}
\setmathfont{Latin Modern Math}
\usepackage{microtype}
\usepackage{enumitem,xurl,needspace}
\usepackage[dvipsnames]{xcolor}
\usepackage{hyperref}
\hypersetup{unicode=true,colorlinks=true,linkcolor=MidnightBlue,urlcolor=MidnightBlue,
 pdftitle={UnsolvedMath Problem 1423},pdfauthor={GPT 6 Astra Ultra},bookmarksnumbered=true}
\usepackage{bookmark,tocloft,titlesec,fancyhdr}
\setlength{\cftsecnumwidth}{4.2em}
\cftsetpnumwidth{2.5em}
\cftsetrmarg{4em}
\titleformat{\section}{\Large\bfseries\raggedright}{\thesection}{0.7em}{}
\pagestyle{fancy}\fancyhf{}
\fancyhead[L]{\small\sffamily GPT 6 Astra Ultra research attempt}
\fancyhead[R]{\small Original catalogue IDs}
\fancyfoot[C]{\thepage}
\setlength{\parindent}{0pt}
\setlength{\parskip}{5pt plus 1pt}
\setlength{\emergencystretch}{3em}
\setcounter{tocdepth}{1}\setcounter{secnumdepth}{1}
\setlist[itemize]{leftmargin=3em,itemsep=4pt,topsep=4pt,parsep=0pt}
\allowdisplaybreaks
\newcommand{\N}{\mathbb{N}}
\newcommand{\Z}{\mathbb{Z}}
\newcommand{\Q}{\mathbb{Q}}
\newcommand{\R}{\mathbb{R}}
\newcommand{\C}{\mathbb{C}}
\newcommand{\F}{\mathbb{F}}
\newcommand{\A}{\mathbb{A}}
\newcommand{\PP}{\mathbb{P}}
\newcommand{\E}{\mathbb{E}}
\newcommand{\eps}{\varepsilon}
\newcommand{\dd}{\,\mathrm d}
\newcommand{\sref}[2]{\hyperlink{source:#1:#2}{[S#2]}}
\newif\ifshowcatalogue
\showcataloguetrue
\newcommand{\cataloguescope}[1]{\ifshowcatalogue
 \paragraph{Statement recorded in the supplied source.}
 #1\fi}
\begin{document}
% UnsolvedMath ID 1423; source catalogue code GRAPH-036

\setcounter{section}{1422}

\section{The Snake-in-the-Box Problem}\label{1423}

{\small\sffamily UnsolvedMath ID 1423 · Graph Theory}

\subsection{Definitions and mathematical statement}

\paragraph{Shared notation.}Unless a section says otherwise, $\N=\{1,2,\ldots\}$,
$\Z$ is the ring of integers, $\Q,\R,\C$ are the rational, real, and complex
fields, $[N]=\{1,\ldots,\lfloor N\rfloor\}$, and $\log$ is the natural
logarithm. For real functions of a parameter tending to infinity,
$f\ll g$ and $f=O(g)$ mean $|f|\leq Cg$ eventually for some fixed $C$;
$f\gg g$ reverses this comparison; $f=o(g)$ means $f/g\to0$;
$f\sim g$ means $f/g\to1$; and $f\asymp g$ means both order bounds.
A subscript on $O$ or $\ll$ specifies allowed constant dependence.
The expression $n^{o(1)}$ in an upper bound means growth smaller than
$n^\epsilon$ for each fixed $\epsilon>0$. Local definitions govern any
exception, finite-field convention, or use of the same symbol for another
object. An asymptotic claim is not automatically an effective algorithm.



The hypercube $Q_n$ has vertex set $\{0,1\}^n$, with an edge between strings differing in exactly one coordinate. An induced path $v_0,\ldots,v_\ell$ has distinct vertices and an edge between two path vertices exactly when their indices differ by one. Its length is $\ell$ edges; its vertex count is $\ell+1$. Determine the maximum length. \sref{1423}{1} \sref{1423}{2}

\paragraph{Statement recorded in the supplied source.}
What is the longest induced path in an n-dimensional hypercube graph? \sref{1423}{1}

\paragraph{Residual task reported by the archive.}

The embedded review (2026-08-17) records: Determine exact values beyond dimension 8 and sharpen asymptotic bounds. \sref{1423}{1}

\subsection{Short English statement}

What is the longest path through a binary hypercube that has no shortcut edges between nonconsecutive path vertices?

\subsection{Research attempt}

\paragraph{Length convention and selected route.}
Write $L(n)$ for the maximum number of edges in an induced path in
$Q_n$. A path with $m$ vertices has $m-1$ edges throughout this attempt.
The route combines a boundary-counting upper bound, a constructive
dimension-lifting lower bound, and a small exhaustive search. Wynn's
primary manuscript [E1] supplies context for more sophisticated snake
and coil constructions; its historical records are not represented as
the current optimal values in every dimension.

\paragraph{An upper bound from edges leaving the path.}
Let $P$ be induced and have $m\ge2$ vertices. The cube is $n$-regular,
and exactly $m-1$ edges have both endpoints in $P$. Summing degrees
of vertices in $P$ therefore counts
\[
                     nm-2(m-1)=(n-2)m+2
\]
edges from $P$ to its complement. Each of the $2^n-m$ outside
vertices receives at most $n$ such edges. Hence for $n\ge2$,
\[
 (n-2)m+2\le n(2^n-m),\qquad
 L(n)\le\left\lfloor\frac{n2^n-2}{2n-2}\right\rfloor-1.
\]
At $n=1$ the single cube edge gives $L(1)=1$ directly; one must not
divide by $2n-2$ there. Inducedness is crucial: if the path vertex
set had extra internal edges, subtracting only $2(m-1)$ from the
degree sum would be wrong.

The bound treats every outside vertex as though it could meet the path
in all $n$ neighbours independently. Those incidences are constrained
by the cube and by the order of the path. Thus the inequality is
generally loose. Equality in it would require every outside vertex
to have all its neighbours inside the path, in addition to the path
being induced. Merely attaining the numerical degree sum does not
construct such a configuration.

\paragraph{Elementary paths and an exponential lifting construction.}
Starting at the all-zero word and flipping coordinates $1,2,\ldots,n$
once each gives a path of length $n$. The Hamming distance between
the vertices after $i$ and $j$ flips is $|i-j|$, so no nonconsecutive
vertices are adjacent. Thus $L(n)\ge n$.

For a stronger recurrent construction, let
$a_0,a_1,\ldots,a_{m-1}$ be any induced path in $Q_n$. In $Q_{n+2}$
take the sequence
\[
 (a_0,00),\ldots,(a_{m-1},00),
 (a_{m-1},01),
 (a_{m-1},11),\ldots,(a_0,11).
\]
Each consecutive pair is adjacent. The two long copies are individually
induced; a vertex in layer $00$ is at distance at least two from a
vertex in layer $11$. The sole vertex in layer $01$ is adjacent to
a vertex of either copy only when its first $n$ coordinates equal
$a_{m-1}$. Thus it introduces precisely its two intended neighbours
and no chord. The new path has $2m+1$ vertices and $2m$ edges, giving
\[
                         L(n+2)\ge2L(n)+2.
\]
For example the exact base $L(1)=1$ yields $L(2r+1)\ge3\cdot2^r-2$,
and the base $L(2)=2$ yields $L(2r)\ge2^{r+1}-2$ for $r\ge1$.
These are valid infinite families, not assertions of sharpness.

\paragraph{An exact exhaustive search with no unsafe symmetry reduction.}
Encode cube vertices by integers $0,\ldots,2^n-1$, with adjacency
when the exclusive-or has exactly one set bit. Translation by a cube
vertex is an automorphism, so every induced path can be translated to
start at zero. This justifies fixing its first vertex. The search then
tries every unused neighbour $w$ of the current endpoint, accepting it
only when it has no neighbour among the earlier path vertices other
than that endpoint. The condition is exactly what is needed to preserve
inducedness.

Inductively, every induced path beginning at zero occurs in this search
tree: its next vertex is among the candidates and passes the condition.
No unproved coordinate-permutation pruning is used. Recording the
largest depth after visiting the whole finite tree therefore proves
the optimum for the searched dimension, not only a lower bound.
The implementation was executed for dimensions one through four and
returned
\[
                            L(1),L(2),L(3),L(4)=1,2,4,7.
\]
The numbers of visited search states were $2,5,22,305$, respectively.
These are complete search-tree counts with the starting vertex fixed.

\paragraph{Concrete witnesses and reproducibility.}
For dimensions three and four the saved integer-encoded witnesses are
\[
                 0,1,3,7,6
                 \quad\hbox{and}\quad
                 0,1,3,7,6,14,12,13.
\]
For example consecutive exclusive-ors in the latter are
$1,2,4,1,8,2,1$, each a power of two; the program also verifies the
nonconsecutive-neighbour exclusion at every extension. The first
sequence has five vertices and four edges, the second eight vertices
and seven edges. The audit files are \texttt{checks\_graphs.py} and
\texttt{results\_graphs.json} in
\path{_Local/Erdos_problems/UnsolvedMath/attempt_audit_2026_10_01/number_theory/}.
No larger-dimensional exhaustive computation is claimed.

\paragraph{Remaining gap and outcome.}
The construction grows roughly like $2^{n/2}$, while the elementary
upper bound is of order $2^{n-1}$. Thus these arguments leave even an
exponential gap in their asymptotic scales. Better concatenations must
control cross-copy chords; better upper bounds must exploit cube
structure beyond the total boundary capacity. The four exact finite
values and the general inequalities do not determine $L(n)$ for all
$n$. \textbf{Outcome: UNRESOLVED.} Neither the construction nor the
small computed values are claimed as new records.

\subsection{Sources}

{\small

\begin{itemize}

\item[\hypertarget{source:1423:1}{S1}] ULAM.AI, \emph{UnsolvedMath}, supplied \texttt{problems.json}, record 1423 (GRAPH-036), statement and background. The embedded literature-review date is 2026-08-17; that is the archive’s review date, not a fresh certification. Dataset landing page: \url{https://huggingface.co/datasets/ulamai/UnsolvedMath}.

\item[\hypertarget{source:1423:2}{S2}] Snake-in-the-box overview. \url{https://en.wikipedia.org/wiki/Snake-in-the-box}. Bibliographic information imported from the supplied record.


\item[E1] Ed Wynn, \emph{Constructing circuit codes
by permuting initial sequences}, arXiv:1201.1647v1 (2012).
\url{https://arxiv.org/abs/1201.1647}. Primary abstract inspected
1 October 2026. It discusses constructions of snakes and coils;
the historical record claims in that abstract are not updated here.

\end{itemize}

}

\label{end:1423}
\end{document}
